Пример работы разработанного программного обеспечения

После загрузки модуля файлы из списков становятся невидимыми и защищёнными.

\begin{lstlisting}[language=bash, frame=single, numbers=left, caption={Добавление файлов в списки}]
artem@artem-pc:~$ echo hidden.txt   > /proc/hidden
artem@artem-pc:~$ echo protected.txt > /proc/protected
\end{lstlisting}

\begin{lstlisting}[language=bash, frame=single, caption={Содержимое директории после загрузки модуля}]
artem@artem-pc:~/OS_CP$ ls
file.txt  protected.txt
\end{lstlisting}

Файл hidden.txt полностью исчез из листинга.

\begin{lstlisting}[language=bash, caption={Отключение режима скрытия}]
artem@artem-pc:~$ echo 1234 > /dev/emblem
\end{lstlisting}

После ввода пароля 1234 файл снова виден:

\begin{lstlisting}[language=bash, caption={Файлы снова видны}]
artem@artem-pc:~/OS_CP$ ls
file.txt  hidden.txt  protected.txt
\end{lstlisting}

Аналогично, попытка удалить или прочитать protected.txt ничего не даёт:

\begin{lstlisting}[language=bash, caption={Защита от удаления и чтения работает}]
artem@artem-pc:~$ rm protected.txt
artem@artem-pc:~$ echo $?
1
artem@artem-pc:~$ cat protected.txt
(nothing appears)
\end{lstlisting}

Отключаем защиту:

\begin{lstlisting}[language=bash, caption={Отключение режима защиты}]
artem@artem-pc:~$ echo 4321 > /dev/emblem
\end{lstlisting}

Теперь операции проходят успешно:

\begin{lstlisting}[language=bash, caption={Файл снова доступен}]
artem@artem-pc:~$ cat protected.txt
This is a protected file
\end{lstlisting}